\chapter*{Introducción}
\addcontentsline{toc}{chapter}{Introducción}
\markboth{Introducción}{Introducción}

\lettrine[lines=2]{L}{os} sistemas estigmáticos son de gran importancia en
la óptica, debido a que resuelven el problema de la formación de una imagen
perfecta de un objeto puntual en la óptica geométrica
\cite{bornwolf1999,Tyc2011Absolute}, con aplicaciones directas en la
microscopía de superresolución y la litografía óptica
\cite{Shen2022,flagello1996}, el procesamiento láser de materiales
\cite{Niziev1999Cutting} y el atrapamiento óptico de partículas
\cite{Kozawa2010Trapping}.

El estigmatismo es un importante punto de partida para el estudio
de la formación de imágenes \cite{singlet2020, bornwolf1999}. Las Superficies Cartesianas aseguran esta
condición en la refracción entre dos medios homogéneos y permiten construir
sistemas mediante su composición \cite{singlet2020,ovoids2023}. Una elección
adecuada de las Superficies Cartesianas permite obtener propiedades como el
aplanatismo y el acromatismo, lo que demuestra su utilidad y potencial en el
diseño óptico \cite{aplan2020,achro2022}.

La minimización de las aberraciones geométricas conserva su importancia aun
en la óptica física, donde estas se manifiestan como aberraciones de fase
que alteran el campo en la imagen \cite{bornwolf1999,goodman2005}. La relación directa con la función aberración de fase y la longitud de camino óptico, permite deducir que en los sistemas estigmaticos esta aberración es nula. Esta propiedad hace de las Superficies Cartesianas un
escenario apropiado para aislar las aberraciones vectoriales, pues permite
estudiar la transformación de amplitud y polarización sin mezclarla con
defectos de camino óptico. Aun así, la difracción produce una distribución
extendida,
descrita por la función de dispersión de punto (PSF, por sus siglas en inglés),
y establece limitaciones de resolución cuya referencia clásica es el límite
de Abbe \cite{Abbe1873,goodman2005}. La descripción escalar no permite
identificar los efectos de la polarización sobre esa distribución. Las
variaciones de amplitud y polarización sobre la pupila pueden deformar la
PSF y reducir el contraste de la imagen aun sin aberración de fase, lo que
motiva su estudio mediante una descripción vectorial
\cite{McGuire1991Response,mcguire1994,flagello1996}.

El impacto de estos efectos alcanza tanto la formación de imágenes como la
interacción de la luz con la materia. En litografía de alta apertura
numérica, la polarización modifica la distribución de intensidad y las
asimetrías de la imagen \cite{flagello1996}. En microscopía de
superresolución, las aberraciones de polarización pueden alterar el mínimo
de intensidad de los haces de agotamiento y los patrones de iluminación
estructurada. En microscopía polarimétrica, pueden introducir errores en
la información de polarización atribuida a la muestra \cite{Shen2022}.
Asimismo, la
polarización condiciona la eficiencia axial y transversal de las pinzas
ópticas \cite{Kozawa2010Trapping} y la absorción de energía que determina la
eficiencia del corte láser \cite{Niziev1999Cutting}. Estas aplicaciones
motivan la caracterización de las alteraciones de amplitud y polarización
introducidas por la refracción, incluso cuando la aberración de fase ya se
ha eliminado.

El tratamiento electromagnético dispone de representaciones integrales
exactas basadas en identidades vectoriales de Green \cite{StrattonChu1939}
y de métodos numéricos como el de elementos finitos (FEM) y el de diferencias
finitas en el dominio del tiempo (FDTD), cuyo costo computacional dificulta
el estudio de sistemas grandes respecto de la longitud de onda
\cite{Shi2019Curved,Wende2022,Song2025}. Las aproximaciones de Debye y
Richards--Wolf permiten un cálculo focal eficiente, pero presentan
restricciones de validez y, en el modelo aplanático ideal de
Richards--Wolf, omiten la transmisión diferencial de Fresnel en las
superficies de la lente
\cite{RichardsWolf1959,Sheppard2013Debye,Shi2019RealLenses,Wang2020Debye}.
Otros métodos incorporan la transmisión local mediante trazado vectorial,
aproximaciones de interfaz plana y operaciones espectrales
\cite{Kim18,Shi2019Curved,Wende2022,Song2025}.\footnote{El presente
desarrollo es independiente de los trabajos de Shi et
al.~\cite{Shi2019Curved} y Song et al.~\cite{Song2025}, conocidos durante
la revisión bibliográfica posterior.} Sin embargo, estos tratamientos
generales no caracterizan las aberraciones vectoriales de las Superficies
Cartesianas, y el antecedente específico de difracción vectorial en estas
superficies considera su apodización, pero excluye los coeficientes de
Fresnel en la comparación focal \cite{GonzalezAcuna2022Oval}. Queda así por
determinar la contribución de la transmisión vectorial a la estructura del
campo focal en ausencia de aberración de fase.

El objetivo de este trabajo es aislar y caracterizar las aberraciones de
origen netamente vectorial que estan presentes aun en ausencia de aberración de
fase. Para ello se calcula el campo eléctrico transmitido hacia el plano focal de una Superficie Cartesiana entre dos medios dieléctricos homogéneos e isótropos, donde el campo incidente es una onda esférica centrada en el punto objeto.
Este sistema es el sistema estigmático más simple, y nos permite examinar los efectos de la refracción vectorial sin la contribución de las aberraciones de camino óptico. 
En este trabajo expresamos la transmisión de forma operacional, dado el campo incidente mediante el operador matricial que se contruyó en este trabajo, damos el campo transmitido sobre la esféra gaussiana centrada en el punto imagen, matriz que a su vez, nos permite caracterizar la aberración vectorial siguiendo los métodos de McGuire y Chipman \cite{McGuire1991Response}. Este campo transmitido, lo propagamos usando un método vectorial que usa la linealidad de la ecuación de Helmholtz para resolver el problema vectorial propagando cada componente lineal como campos escalares, respetando la condición de transversalidad \cite{Song2025}. Bajo la aproximación de Fresnel el campo en el plano focal queda determinado por una transformación de Fourier de la aliación del operador de transmisión sobre el campo incidente. Resultado que se reduce a transformadas de Hankel cuando la iluminación incidente presenta simetría axial, formulación que permite interpretar con mayor facilidad los efectos vectoriales sobre el plano focal.  Nuestro formalismo tiene en cuenta los efectos vectoriales tales como el acoplamiento y modulación de las polarizaciones, y la componente longitudinal del campo. Este análisis de las aberraciones vectoriales en sistemas estigmáticos es un primer paso hacia su posterior corrección. La composición de las operaciones permite, además, extender el cálculo a sistemas formados por varias Superficies Cartesianas.
